% \iffalse meta-comment
%%
%% shmuelsemtex.dtx
%%
%% Copyright (C) 2026 Seymour J. Metz
%%
%% This material is subject to the MIT License.
%% See the LICENSE file in the root of this distribution, or
%% <https://github.com/shmuelmetz/LaTeX-Semantic-Markup>.
%%
%% This is the shmuelsemtex package, part of the
%% LaTeX-Semantic-Markup distribution.
%%
%<*driver>
\documentclass[full]{l3doc}
\usepackage{mathrsfs}
\usepackage{bm}
\usepackage{amssymb}
\usepackage{fontspec}
\newfontfamily\hebrewfont{David CLM}[Script=Hebrew]
\TeXXeTstate=1
\setlength\emergencystretch{3em}
\usepackage{cleveref}
% "<cref> (<name>) on page <n>"
\NewDocumentCommand\pagecrefn{m}{\cref{#1} (\nameref{#1}) on \cpageref{#1}}
\begin{document}
  \DocInput{shmuelsemtex.dtx}
\end{document}
%</driver>
% \fi
%
% \title{The \pkg{shmuelsemtex} package}
% \author{Seymour J. Metz (H: Shmuel {\hebrewfont\beginR שמואל בן לייביש ולאה\endR}) \\ \url{https://mason.gmu.edu/~smetz3} \\ arXiv: \url{http://arxiv.org/a/metz_s_1} \\ Stack Exchange: \url{https://stackexchange.com/users/8873368/shmuel} \\ \TeX{} Stack Exchange: \url{https://tex.stackexchange.com/users/110591/shmuel}}
% \date{2026-09-01}
%
% \maketitle
%
% \begin{documentation}
%
% \section{Purpose}
% \label{sec:purpose}
%
% This package provides macros that record the \emph{semantic role}
% of a piece of mathematical text --- ``this is the display name of a
% category,'' ``this is a monomorphism arrow'' --- separately from
% how that role is \emph{typeset}. An author writing
% \cs{catName}\Arg{C} is not choosing a font; they are asserting that
% \texttt{C} names a category, and the package decides what a
% category name looks like.
%
% The category-theory naming layer above (\cs{catName}, \cs{morphName},
% \cs{morphMono}\slash \cs{morphEpi}, etc.) is only part of the package.
% Quantifiers (\cs{equant}\slash \cs{uquant}) and set-builder notation
% (\cs{set}\slash \cs{union}\slash \cs{intersection}), documented in
% \pagecrefn{sec:quant-set} below, are general first-order-logic and
% set-theory constructs, not specific to category theory, and are
% equally usable in any mathematical document.
%
% \begin{quote}
% \textbf{Design goal:} a paper submitted to different journals should
% not need different markup for their different house styles, only
% different setup. Concretely:
% \cs{catName}\Arg{Set} (or any semantic macro, anywhere in the
% document body) must be byte-for-byte identical whether the target
% is arXiv, journal~X's house style, or journal~Y's house style. Only
% \emph{setup} may differ: the \texttt{style} package option and the
% setup commands \cs{setupsemanticmarkup} (see
% \pagecrefn{sec:style-switching}), \cs{setupquant} and \cs{setupset} (see
% \pagecrefn{sec:quant-set}) --- never the markup at the use sites.
%
% This is not just an API-ergonomics nicety; it is the whole point of
% the package, and it generalizes the ``macros do the heavy lifting''
% principle used throughout the design decisions below into a
% concrete, checkable test:
% \textbf{the journal-portability test} --- \emph{if satisfying a
% house-style change would ever require touching the markup at the
% use sites (as opposed to setup), the design is wrong.} Every design
% question in this package, past and future, is judged against this
% test.
% \end{quote}
%
% Two consequences follow directly from the test above, and both are
% enforced throughout this package:
% \begin{enumerate}
% \item Two things that look similar on the page but mean different
%   things mathematically must not share a macro, even if that means
%   near-identical default typesetting for two commands (otherwise a
%   house style that \emph{does} want them to look different could
%   not ask for that without editing the body).
% \item One concept must not be given two names merely because it is
%   styled differently in different contexts (otherwise the author
%   would have to choose between the two names based on a
%   presentation decision, rather than on what they mean) --- and no
%   macro's presentation may be hard-coded at all: every macro's
%   actual rendering is decided by the active \emph{style}
%   (\pagecrefn{sec:style-switching}), never per use.
% \end{enumerate}
%
% \section{Background}
% \label{sec:background}
%
% The notation in this package was first written as private macros in
% two papers by its author, \emph{Local Coordinate Spaces: a proposed
% unification of manifolds with fiber bundles, and associated machinery}
% \cite{lcs} and \emph{Towards a taxonomy of atlases and of morphisms
% between them} \cite{matlas}. Those macros were refactored into this
% package, which the papers now load; \pagecrefn{sec:history} gives the
% details.
%
% \section{Internal naming}
%
% Package-internal \pkg{expl3} functions use the module prefix
% \texttt{lsm} (for \LaTeX\ Semantic Markup), per the naming
% convention \texttt{\textbackslash lsm\_}\emph{module}\texttt{\_}%
% \emph{description}\texttt{:}\emph{args}. These are not intended for
% direct use in documents.
%
% \section{Style switching}
% \label{sec:style-switching}
%
% Every semantic macro in this package (\cs{catName}, \cs{morphMono},
% \ldots) is a thin public wrapper around an internal \pkg{expl3}
% function of the same shape (\cs{lsm_cat_name:n}, \cs{lsm_morph_mono:%
% }, \ldots) that does the actual typesetting. Those internal
% functions are not fixed: each one is a \emph{presentation hook},
% (re)definable as a group, all at once, by a named \emph{style}. This
% hook architecture is not new machinery bolted on for style
% switching; it already existed, one hook per macro, from the first
% version of this package (motivated by the same
% typography/semantics separation), and style switching is built as a
% dispatch layer on top of exactly those existing hooks, rather than
% as some separate mechanism.
%
% A style is normally selected as a package option:
% \begin{verbatim}
% \usepackage[style=default]{shmuelsemtex}   % built in, default
% \usepackage[style=plain]{shmuelsemtex}     % built in, alternate
% \end{verbatim}
% and can also be changed anywhere in the document with
% \cs{setupsemanticmarkup} (below). Quantifier and set-builder
% rendering have their own setup commands, \cs{setupquant} and
% \cs{setupset} (\pagecrefn{sec:quant-set}). Either way, the markup at the
% use sites --- every \cs{catName}, every \cs{morphCompose},
% everything --- is written once and does not vary with the setup;
% moving between, say, an arXiv upload and a journal submission
% changes only the setup. This is the direct implementation of the
% journal-portability test in \pagecrefn{sec:purpose}.
%
% The mechanism is \pkg{l3keys2e} package-option processing (the
% standard modern \pkg{expl3} idiom for \texttt{key=value}
% \cs{usepackage} options; the same mechanism many CTAN packages use
% for exactly this ``pick a coherent bundle of choices at once''
% pattern), with a single \texttt{.choice:} key, \texttt{style}, in
% the key family \texttt{lsm}. Each choice's code is one function
% (\cs{lsm_style_default:}, \cs{lsm_style_plain:}) that reassigns
% \emph{every} presentation hook in the package in one place, so a
% style is always a complete, coherent bundle --- never a partial
% patch the author has to assemble by hand from individual macro
% redefinitions. \texttt{style} defaults to \texttt{default} via
% \texttt{.initial:n}, which runs the named choice's code at
% \cs{keys_define:nn} time, so the package works with no options at
% all; any explicit \texttt{style=...} package
% option is applied afterward via \cs{ProcessKeysPackageOptions}.
%
% Only two styles are built in for now, both illustrating the
% mechanism rather than encoding any real journal's house style yet:
% \begin{description}
% \item[\texttt{default}] The presentation choices already justified,
%   macro by macro, in the sections below (script/bold-upright
%   auto-detected category names, bold sequence names, tail/twohead
%   arrows for mono/epi, \ldots).
% \item[\texttt{plain}] Every \texttt{*Name} family macro
%   (\cs{semName}, \cs{catName}, \cs{catSeqName}, \cs{objName},
%   \cs{morphName}, \cs{morphSeqName}, \cs{chartName}, \cs{atlasName},
%   \cs{coordName}) renders completely unstyled, ambient math
%   typeface --- no script, no bold, no auto-detection. The
%   mono/epi arrows and the compose/compose-head/compose-tail glyphs
%   are left as in \texttt{default}: those are closer to fixed
%   mathematical convention (an arrow shape for a monomorphism,
%   distinct glyphs for three distinct composition operations) than
%   to journal-specific name typography, so a demonstration ``plain''
%   style does not touch them --- a real house-style-specific style
%   built later might, but that is a further design question, not
%   answered here.
% \end{description}
% Where no predefined style matches the house style, a new style is
% added as one more \texttt{style / \meta{name} .code:n} choice and
% one more \cs{lsm_style_\meta{name}:} function that reassigns the
% same complete set of hooks; it never means touching the public
% macros or the markup at the use sites. The package does not yet
% provide a user-level command for defining a style in a document's
% own preamble.
%
% \begin{quote}
% \textbf{Runtime setup.} \cs{setupquant} and \cs{setupset}
% (\pagecrefn{sec:quant-set}) are runtime setup commands built on plain
% \pkg{l3keys}, called once, away from every use site:
% \begin{verbatim}
% \setupquant {parentheses=multiple,subscript=stacked}
% \end{verbatim}
% after which every \cs{equant}\slash \cs{uquant} call, anywhere in the
% document, renders accordingly, unchanged at the call site. Unlike a
% package option, a runtime command can be invoked anywhere a document
% needs it, not only on the \cs{usepackage} line.
%
% The \texttt{style} key is available at runtime in the same way, via
% \cs{setupsemanticmarkup}, a thin wrapper around \cs{keys_set:nn} on
% the key family defined for the package options:
% \begin{verbatim}
% \setupsemanticmarkup {style=plain}
% \end{verbatim}
% Both forms are kept: \cs{keys_set:nn} is what
% \cs{ProcessKeysPackageOptions} calls to apply a package option, and
% \cs{setupsemanticmarkup} makes the same call available at any point
% in the body.
%
% This package's styles remain one coherent \texttt{style} bundle
% covering every macro at once, unlike \cs{setupquant}\slash \cs{setupset}'s
% split into independent per-concept-family setup commands: a
% journal's house style is ordinarily one holistic choice applied
% consistently everywhere, not an independent axis per semantic
% family the way quantifier-rendering-taste and set-builder-rendering-%
% taste are independent of each other for a single author. If a real
% future style needs finer-grained, per-family control, splitting
% \texttt{style} into per-family keys (mirroring \texttt{quant}/%
% \texttt{set} exactly) is a natural extension of this same
% \pkg{l3keys} foundation, not a different mechanism.
% \end{quote}
%
% \section{Abbreviations used in macro names}
%
% Every abbreviated word appearing in a user-level macro name is
% spelled out here, once, so no reader has to guess it from context.
% Family and role prefixes that are already complete English words
% (\texttt{chart}, \texttt{atlas}, \texttt{name}, \texttt{compose})
% are not abbreviations and are not listed.
%
% \begin{description}
% \item[\texttt{sem}] Semantic
% \item[\texttt{cat}] Category
% \item[\texttt{obj}] Object
% \item[\texttt{morph}] Morphism
% \item[\texttt{coord}] Coordinate
% \item[\texttt{seq}] Sequence (as in \cs{catSeqName}: a
%   \emph{sequence of} the preceding concept, not that concept itself)
% \item[\texttt{mono}] Monomorphism
% \item[\texttt{epi}] Epimorphism
% \end{description}
%
% \section{User-level command families}
%
% User-level commands are grouped by the semantic domain they tag,
% each using its own name prefix:
% \begin{itemize}
%   \item \texttt{sem} (Semantic) --- generic / not-yet-classified
%     semantic text
%   \item \texttt{cat} (Category) --- categories (in the
%     category-theoretic sense)
%   \item \texttt{obj} (Object) --- objects of a category
%   \item \texttt{morph} (Morphism) --- morphisms between objects
%   \item \texttt{chart} --- charts
%   \item \texttt{atlas} --- atlases of charts
%   \item \texttt{coord} (Coordinate) --- coordinate systems/functions
% \end{itemize}
% What follows documents what each macro \emph{means} and why it is
% (or is not) distinct from its neighbors --- that does not change
% with the active style. What each macro actually renders as, under
% each built-in style, is in \pagecrefn{sec:style-switching} and the
% Implementation section below.
%
% \subsection{Generic}
%
% \begin{function}{\semName}
% \begin{syntax}
% \cs{semName}\marg{text}
% \end{syntax}
% \cs{semName}\Arg{text} (Sem = Semantic) tags \texttt{text} as
% semantically significant without committing to any of the more
% specific domains below. Use this only as a placeholder until a
% more specific macro applies, or for genuinely domain-agnostic
% labels.
% \end{function}
%
% \subsection{Categories}
%
% \begin{function}{\catName}
% \begin{syntax}
% \cs{catName}\marg{text}
% \end{syntax}
% \cs{catName}\Arg{text} (Cat = Category) is the display name of a
% single category, whether an arbitrary/variable category symbol
% (e.g.\ \cs{catName}\Arg{C}) or a long, conventionally-fixed name
% (e.g.\ \cs{catName}\Arg{Set}). The author writes the same macro
% either way and never separately decides a typeface: see the
% discussion below of symbols and names for why this is
% one auto-detecting macro rather than two.
% \end{function}
%
% \begin{quote}
% \textbf{One macro for symbols and names:} a category may be written
% with an arbitrary symbol, such as \texttt{C}, or with a name, such as
% \texttt{Set} or \texttt{Top}. The difference in typeface is a
% convention about the length of the name, with no semantic distinction
% behind it. Two macros would force the author to decide a typographic
% question as a separate step from expressing what they mean.
% \cs{catName} is therefore a single macro that decides for itself: the
% author writes \cs{catName}\Arg{Set} or \cs{catName}\Arg{C}
% identically and never chooses upright vs.\ script. Under the
% \texttt{default} style an arbitrary symbol renders in script and a
% name renders bold upright, by the test in
% \pagecrefn{sec:cat-symbols}; see \pagecrefn{sec:style-switching} for
% how a different style swaps that choice as a unit, without the author
% touching \cs{catName}\Arg{Set} itself.
% \end{quote}
%
% \begin{quote}
% \label{sec:cat-symbols}
% \textbf{Symbols vs.\ names:} token count alone does not separate
% symbols from names. \texttt{C\^{}2}, \texttt{E\^{}i} and
% \texttt{XY} (a product-space category \(X \times Y\)) are arbitrary
% symbols, yet none is a single token.
% What distinguishes them from \texttt{Set}/\texttt{Top} is whether
% the argument reads as a \emph{decorated or concatenated symbol} (a
% base letter carrying a sub/superscript, or two symbol-like letters
% standing for a product) versus a genuine multi-letter \emph{word}.
% \cs{catName} tests, in order: (1) does the argument contain \texttt{\^{}} or \texttt{\_{}}
% (a sub/superscript decoration) --- if so, script, regardless of what
% the decoration itself contains (this is why \texttt{E\^{}i}'s
% lower-case \texttt{i} never has to be examined: the whole argument
% is already settled as ``decorated symbol'' by the presence of
% \texttt{\^{}} alone); (2) is it a single token --- if so, script;
% (3) is every letter in the argument already upper case ---
% if so, script (this is what catches \texttt{XY}: two upper-case
% letters standing for a product, not a lower-case-bearing word); (4)
% otherwise, bold upright.
% \end{quote}
%
% \begin{function}{\catSeqName}
% \begin{syntax}
% \cs{catSeqName}\marg{text}
% \end{syntax}
% \cs{catSeqName}\Arg{text} (Cat = Category, Seq = Sequence) is the
% display name of a \emph{sequence or family of categories}, e.g.
% $(\mathscr{A}_\alpha)_{\alpha \in A}$. This is deliberately a
% distinct command from \cs{catName}, not a styling variant of it: a
% sequence of categories is not itself a category, and collapsing the
% two into one macro differentiated only by an optional ``bold''
% switch would silently re-introduce the typography/semantics
% conflation this package exists to remove.
% \end{function}
%
% \subsection{Objects}
%
% \begin{function}{\objName}
% \begin{syntax}
% \cs{objName}\marg{text}
% \end{syntax}
% \cs{objName}\Arg{text} (Obj = Object) is the display name of an
% object of a category.
% Under the \texttt{default} style the presentation is unstyled
% (objects are typically bare math-mode variables in practice); the
% macro exists so that document text records the semantic role even
% when the current style's rendering does nothing visible.
% \end{function}
%
% \subsection{Morphisms}
%
% \begin{function}{\morphName}
% \begin{syntax}
% \cs{morphName}\marg{text}
% \end{syntax}
% \cs{morphName}\Arg{text} (Morph = Morphism) is the display
% name/label of a morphism (e.g.\ a function symbol).
% \end{function}
%
% \begin{quote}
% \textbf{A morphism vs.\ the class of all morphisms:} by
% the standard definition, a single morphism always has a specific
% source and target --- it is an element of a particular hom-set.
% That is a different notion from ``the class of all morphisms/arrows
% in a category,'' which has no reference to specific endpoints.
% \cs{morphName} (a label for one morphism, e.g.\ $f$) is already
% built around the correct, standard reading.
%
% \texttt{morphAll} (all morphisms of a category) and
% \texttt{morphBetween} (morphisms from one specific object to
% another, a hom-set) are natural names
% for these two notions but are not built here, pending a decision on
% naming.
% \end{quote}
%
% \begin{quote}
% \textbf{Function symbols vs.\ names:} \cs{morphName}\slash
% \cs{morphSeqName} use the same test as \cs{catName}
% (\pagecrefn{sec:cat-symbols}): decoration (\texttt{\^{}}/%
% \texttt{\_{}}) or single token or all-upper-case letters -> italic
% (or bold italic); otherwise -> upright (or bold upright). Single
% tokens include single \emph{control sequences} such as \cs{pi} and
% \cs{rho}, whatever the length of their names. As with
% \texttt{C\^{}2}, \texttt{f\^{}2} is an arbitrary symbol though not
% a single token: a decorated base symbol is still an arbitrary symbol,
% not a name.
% \end{quote}
%
% \begin{function}{\morphSeqName}
% \begin{syntax}
% \cs{morphSeqName}\marg{text}
% \end{syntax}
% \cs{morphSeqName}\Arg{text} (Morph = Morphism, Seq = Sequence) is
% the display name of a \emph{sequence or tuple of morphisms/%
% functions}, e.g.\ $\bm{f} \mathrel{:=} (f_1, f_2)$. Like
% \cs{catSeqName}
% vs.\ \cs{catName}, this is a genuinely distinct \pkg{expl3} function
% from \cs{morphName} (\cs{lsm_morph_seqname:n}, not
% \cs{lsm_morph_name:n} with a flag): a sequence of functions is not
% itself a function. Auto-detects the same way \cs{morphName} does,
% but bold throughout --- see the note on function symbols above.
% See the discussion
% below.
% \end{function}
%
% \begin{quote}
% \textbf{morphSeqName vs.\ morphName:} a sequence of
% functions is not a function, the same principle as the
% \cs{catSeqName}\slash \cs{catName} and compose/compose-head/compose-tail
% distinctions in this section. \cs{morphSeqName} is
% built accordingly as a genuinely
% distinct \pkg{expl3} function, \cs{lsm_morph_seqname:n}, not
% \cs{lsm_morph_name:n} called with a different flag.
% \end{quote}
%
% \begin{function}{\morphMono}
% \begin{syntax}
% \cs{morphMono}
% \end{syntax}
% \cs{morphMono} (Mono = Monomorphism) is the arrow symbol for a
% monomorphism.
% \end{function}
%
% \begin{quote}
% \textbf{Tail, not hook:} a hook ($A \hookrightarrow B$) represents
% an inclusion map. A tail ($A \rightarrowtail B$) represents a
% monomorphism. These are two distinct notions with two distinct
% symbols (every inclusion is a monomorphism, but not every
% monomorphism is an inclusion); conflating them under one glyph is
% exactly the typography/semantics conflation this package exists to
% prevent. \cs{morphMono} therefore renders as \cs{rightarrowtail}
% (↣, from \pkg{amssymb}; see \pagecrefn{sec:font-guard}), matching
% the category-theory
% convention that draws the same hook-vs-tail distinction.
% The hook arrow is \cs{morphInclArrow}.
% \end{quote}
%
% \begin{function}{\morphEpi}
% \begin{syntax}
% \cs{morphEpi}
% \end{syntax}
% \cs{morphEpi} (Epi = Epimorphism) is the arrow symbol for an
% epimorphism.
% This is kept distinct from \cs{morphMono} because mono/epi is a
% real mathematical distinction (injective vs.\ surjective in the
% relevant sense), not a choice of arrowhead glyph.
% \end{function}
%
% \begin{function}{\morphCompose}
% \begin{syntax}
% \cs{morphCompose}\oarg{label}
% \end{syntax}
% \cs{morphCompose} is the general morphism-composition operator
% (default $\circ$). An optional argument
% \cs{morphCompose}\oarg{label} places \texttt{label} over the
% operator.
% \end{function}
%
% \begin{function}{\morphComposeHead}
% \begin{syntax}
% \cs{morphComposeHead}\oarg{label}
% \end{syntax}
% \cs{morphComposeHead} (default $\odot$) is composition taken from
% the head of a sequence of functions. An optional argument
% \cs{morphComposeHead}\oarg{label} places \texttt{label} over the
% operator.
% \end{function}
%
% \begin{function}{\morphComposeTail}
% \begin{syntax}
% \cs{morphComposeTail}\oarg{label}
% \end{syntax}
% \cs{morphComposeTail} (default $\cdot$) is composition taken from
% the tail of a sequence of functions. An optional argument
% \cs{morphComposeTail}\oarg{label} places \texttt{label} over the
% operator.
% \end{function}
%
% \begin{quote}
% \textbf{Three distinct compositions:} plain composition, ``compose head,'' and ``compose
% tail'' are genuinely different mathematical operations, not
% glyph/styling variants of one operation. They are built accordingly
% as three separate \pkg{expl3} functions ---
% \cs{lsm_morph_compose:n}, \cs{lsm_morph_compose_head:n},
% \cs{lsm_morph_compose_tail:n} --- each with its own presentation
% hook, the same treatment already given to \cs{catName}/
% \cs{catSeqName} (\cs{lsm_cat_name:n}\slash \cs{lsm_cat_seqname:n}): one
% function with a flag would have re-introduced the
% typography/semantics conflation this package exists to remove.
% \end{quote}
%
% \begin{function}{\morphId}
% \begin{syntax}
% \cs{morphId}\oarg{object}
% \end{syntax}
% \cs{morphId} (Id = Identity) is the identity morphism. The optional
% \meta{object} is the object it acts on, written as a subscript:
% $\operatorname{Id}_A$. Without it, the bare $\operatorname{Id}$ is
% written.
% \end{function}
%
% \begin{function}{\morphSeqId}
% \begin{syntax}
% \cs{morphSeqId}\oarg{object}
% \end{syntax}
% \cs{morphSeqId} is a sequence or tuple of identity morphisms, in
% bold: $\operatorname{\mathbf{Id}}_A$. The optional \meta{object} is
% written as in \cs{morphId}.
% \end{function}
%
% \begin{function}{\morphIncl}
% \begin{syntax}
% \cs{morphIncl}\oarg{included,containing}
% \end{syntax}
% \cs{morphIncl} (Incl = Inclusion) is the inclusion map of
% \meta{included} into \meta{containing}, written with the included
% object as a superscript and the containing object as a subscript:
% $\operatorname{Id}^A_B$. The optional argument is a comma list of at
% most two items; more than two is an error. With one item, only the
% superscript is written; with none, the bare $\operatorname{Id}$. An
% object containing a comma goes in braces. Inclusion is a separate
% command from \cs{morphId}, even where the two render alike, so that
% a style can tell them apart.
% \end{function}
%
% \begin{function}{\morphSeqIncl}
% \begin{syntax}
% \cs{morphSeqIncl}\oarg{included,containing}
% \end{syntax}
% \cs{morphSeqIncl} is a sequence or tuple of inclusion maps, in bold:
% $\operatorname{\mathbf{Id}}^A_B$. The optional argument is as in
% \cs{morphIncl}.
% \end{function}
%
% \begin{function}{\morphInclArrow}
% \begin{syntax}
% \cs{morphInclArrow}
% \end{syntax}
% \cs{morphInclArrow} is the arrow symbol for an inclusion map,
% $\hookrightarrow$, distinct from the tail arrow of \cs{morphMono}.
% \end{function}
%
% \subsection{Charts, atlases, and coordinates}
%
% \begin{function}{\chartName}
% \begin{syntax}
% \cs{chartName}\marg{text}
% \end{syntax}
% \cs{chartName}\Arg{text} is the display name of a chart: a map
% relating part of a space to a model space.
% \end{function}
%
% \begin{function}{\atlasName}
% \begin{syntax}
% \cs{atlasName}\marg{text}
% \end{syntax}
% \cs{atlasName}\Arg{text} is the display name of an atlas: a
% collection of charts.
% \end{function}
%
% \begin{function}{\coordName}
% \begin{syntax}
% \cs{coordName}\marg{text}
% \end{syntax}
% \cs{coordName}\Arg{text} (Coord = Coordinate) is the display name
% of a coordinate system or coordinate function.
% \end{function}
%
% \begin{quote}
% \textbf{Placeholder:} these three commands are
% provided as minimal, unstyled hooks under the \texttt{default} style
% so the semantic distinction exists in document markup from the
% start; being ordinary presentation hooks like every other macro in
% this package, they already participate correctly in style switching
% (\pagecrefn{sec:style-switching}) with no further work needed on that
% front. Whether ``coordinate'' needs sub-kinds (e.g.\ a coordinate
% \emph{function} vs.\ a coordinate \emph{value} vs.\ a coordinate
% \emph{index}) is a separate, still-open semantic question and should
% be settled against real usage before this cluster is trusted as
% final.
% \end{quote}
%
% \subsection{Setup}
% \label{sec:setup}
%
% \begin{function}{style}
% \begin{syntax}
% \cs{usepackage}[style=\meta{name}]\{shmuelsemtex\}
% \end{syntax}
% Selects the style, the complete set of presentation choices for
% every macro in \pkg{shmuelsemtex}. The built-in styles are
% \texttt{default}, used when no option is given, and \texttt{plain};
% \pagecrefn{sec:style-switching} describes both.
% \end{function}
%
% \begin{function}{\setupsemanticmarkup}
% \begin{syntax}
% \cs{setupsemanticmarkup}\marg{style=\meta{name}}
% \end{syntax}
% Runtime counterpart to the \texttt{style} package option: sets the
% same key, and can be called anywhere in the document, not only on
% the \cs{usepackage} line.
% \end{function}
%
% \subsection{Quantifiers}
% \label{sec:quant-set}
%
% \begin{function}{\equant, \uquant}
% \begin{syntax}
% \cs{equant}\marg{variables}\marg{formula}
% \cs{uquant}\marg{variables}\marg{formula}
% \end{syntax}
% Existential (\cs{equant}) and universal (\cs{uquant}) quantification
% of \meta{formula} over \meta{variables}, a comma list of bound
% variables. How the bound variables are placed and parenthesized is
% set with \cs{setupquant}.
% \end{function}
%
% \begin{function}{\setupquant}
% \begin{syntax}
% \cs{setupquant}\marg{key=value list}
% \end{syntax}
% Sets the rendering of \cs{equant} and \cs{uquant}, independently of
% the style. The keys:
% \begin{description}
% \item[\texttt{subscript}] \texttt{none} (initial): the quantifier
%   symbol followed by the variables, joined by the separator;
%   \texttt{stacked}: the variables stacked under an enlarged
%   quantifier symbol; \texttt{multiple} (the value when the key is
%   given alone): one quantifier symbol per variable, each with its
%   variable as a subscript.
% \item[\texttt{parentheses}] \texttt{none} (initial), \texttt{single}
%   or \texttt{multiple} (the value when the key is given alone). Any
%   value other than \texttt{none} puts the quantifier in parentheses;
%   with \texttt{subscript=multiple}, each quantifier separately.
% \item[\texttt{separator}] The token list written between the
%   variables and before the formula: empty initially, \texttt{.} when
%   the key is given alone.
% \end{description}
% \end{function}
%
% \subsection{Set-builder notation}
%
% \begin{function}{\set}
% \begin{syntax}
% \cs{set}\marg{elements}\oarg{propositions}
% \cs{set}\marg{elements}\oarg{propositions}*
% \end{syntax}
% Set-builder notation: \meta{elements} in braces, followed, when
% \meta{propositions} is given, by the separator and the
% propositions. \meta{propositions} is a comma list; each comma is
% typeset as $\land$, so \texttt{\textbackslash set\{x\}[p,q]}
% renders as $\{x \mid p \land q\}$. The starred form
% splits the line between the elements and the propositions, for use
% inside a \env{multline} environment, since \LaTeX{} cannot otherwise
% break a line inside a macro argument. The star follows the
% arguments.
% \end{function}
%
% \begin{function}{\setupset}
% \begin{syntax}
% \cs{setupset}\marg{key=value list}
% \end{syntax}
% Sets the rendering of \cs{set}, \cs{union} and \cs{intersection}.
% The keys:
% \begin{description}
% \item[\texttt{separator}] The symbol between the elements and the
%   propositions of \cs{set}; initially $\mid$.
% \item[\texttt{subscript}] Placement of the indices of \cs{union} and
%   \cs{intersection}: \texttt{stacked} (initial), stacked under an
%   enlarged operator; \texttt{multiple}, a comma-separated subscript.
% \end{description}
% \end{function}
%
% \subsection{Indexed union and intersection}
%
% \begin{function}{\union, \intersection}
% \begin{syntax}
% \cs{union}\oarg{indices}\marg{set}
% \cs{intersection}\oarg{indices}\marg{set}
% \end{syntax}
% Indexed union ($\bigcup$) and intersection ($\bigcap$) of
% \meta{set} over \meta{indices}, a comma list. Without
% \meta{indices}, the bare operator precedes \meta{set}. Index
% placement is set with \cs{setupset}'s \texttt{subscript} key.
% \end{function}
%
% \subsection{Sequence names}
%
% \begin{function}{\seqname}
% \begin{syntax}
% \cs{seqname}\marg{text}
% \end{syntax}
% The name of a sequence, set or tuple: bold upright for a
% multi-character name, bold italic for a single character, the same
% length-based split \cs{catName} uses. It is a separate macro because
% a sequence is a different kind of object from a category.
% \end{function}
%
% \section{History}
% \label{sec:history}
%
% The notation in this package was first written as private macros in
% two papers by the package's author, \emph{Local Coordinate Spaces: a
% proposed unification of manifolds with fiber bundles, and associated
% machinery} \cite{lcs} and \emph{Towards a taxonomy of atlases and of
% morphisms between them} \cite{matlas}, which share one set of
% definitions. The papers' macros for categories, functions,
% composition and arrows named their typography, or named a construct
% only loosely; this package gives the same constructs names that say
% what they are, and moves their presentation into styles. The
% quantifier, set-builder, union and intersection and
% sequence-name commands were already semantic, and keep the papers'
% names: in the papers they form a block marked as belonging in a
% package, copyright 2016, whose header describes them as semantic
% markup for common mathematical constructs and grants permission to
% incorporate them, with attribution, in any freely redistributable
% \LaTeX{} package. The sources of the papers now load this package; the versions
% on arXiv use the original macros. The papers' underset-variant
% commands (\texttt{\textbackslash full},
% \texttt{\textbackslash maxfull}, \texttt{\textbackslash maximal},
% \texttt{\textbackslash strict}, \texttt{\textbackslash Triv}) are
% specific to those papers and remain in them.
%
% \begin{center}
% \begin{tabular}{ll}
% In the papers & In \pkg{shmuelsemtex} \\
% \hline
% \texttt{\textbackslash catname}, \texttt{\textbackslash cat} & \cs{catName} \\
% \texttt{\textbackslash catseqname} & \cs{catSeqName} \\
% \texttt{\textbackslash funcname} & \cs{morphName} \\
% \texttt{\textbackslash funcseqname} & \cs{morphSeqName} \\
% \texttt{\textbackslash compose} & \cs{morphCompose} \\
% \texttt{\textbackslash composeh} & \cs{morphComposeHead} \\
% \texttt{\textbackslash composet} & \cs{morphComposeTail} \\
% \texttt{\textbackslash into} & \cs{morphMono} \\
% \texttt{\textbackslash onto} & \cs{morphEpi} \\
% \texttt{separater} key of \texttt{\textbackslash setupquant}
%   and \texttt{\textbackslash setupset} & \texttt{separator} \\
% \hline
% \texttt{\textbackslash equant}, \texttt{\textbackslash uquant},
%   \texttt{\textbackslash setupquant} & unchanged \\
% \texttt{\textbackslash set}, \texttt{\textbackslash setupset},
%   \texttt{\textbackslash union}, \texttt{\textbackslash intersection}
%   & unchanged \\
% \texttt{\textbackslash seqname} & unchanged \\
% \end{tabular}
% \end{center}
%
% \begin{thebibliography}{9}
% \bibitem{lcs} Seymour J. Metz, \emph{Local Coordinate Spaces: a
%   proposed unification of manifolds and fiber bundles, and associated
%   machinery}, arXiv:1801.05775v1 [math.AT], 17 January 2018.
%   \url{https://arxiv.org/abs/1801.05775}
% \bibitem{matlas} Seymour J. Metz, \emph{Towards a taxonomy of atlases
%   and of morphisms between them}, arXiv:1906.11690v1 [math.AT], 25 June
%   2019. \url{https://arxiv.org/abs/1906.11690}
% \end{thebibliography}
%
% \end{documentation}
%
% \begin{implementation}
%
% \section{Implementation}
%
%<*package>
%    \begin{macrocode}
\ProvidesExplPackage{shmuelsemtex}{2026-09-01}{0.6.0}
  {Semantic markup macros for LaTeX (LSM)}
\RequirePackage{xparse}
\RequirePackage{amsmath}
\RequirePackage{bm}
\RequirePackage{scalerel}
\RequirePackage{l3keys2e}
%    \end{macrocode}
%
% \subsection{Font packages: load only if genuinely needed}
% \label{sec:font-guard}
%
% \cs{mathscr} (needed by \cs{catName}\slash \cs{catSeqName}) and
% \cs{twoheadrightarrow}\slash \cs{rightarrowtail} (needed by \cs{morphEpi}/
% \cs{morphMono} respectively) are what this package actually uses
% from \pkg{mathrsfs}/\pkg{amssymb}. Checking either one of
% \cs{twoheadrightarrow}\slash \cs{rightarrowtail} is enough to guard both,
% since both are amssymb-only symbols (neither is defined by
% \pkg{amsmath} alone) that any font package providing one of them is
% overwhelmingly likely to also provide the other. Requiring both
% packages unconditionally breaks documents that already use a
% comprehensive math font package: \pkg{stix2}, for example, already
% defines both \cs{mathscr} and \cs{twoheadrightarrow}, and loading
% \pkg{amsfonts} (pulled in by \pkg{amssymb}) on top of it exceeds
% LaTeX's limit on symbol fonts (``Too many symbol fonts declared'').
% Guarded with \cs{ifdefined} (a plain \TeX{}
% primitive, unaffected by this file's own \pkg{expl3} catcode regime)
% so a paper that already has an equivalent font package never pays
% for these two a second time.
%    \begin{macrocode}
\ifdefined\mathscr\else\RequirePackage{mathrsfs}\fi
\ifdefined\twoheadrightarrow\else\RequirePackage{amssymb}\fi
%    \end{macrocode}
%
% \subsection{Internal variables}
%
% \begin{variable}[no-user-doc]{\l_lsm_str_tl, \l_lsm_upper_tl}
% Scratch variables used by \cs{lsm_if_arbitrary_symbol:nTF} (below)
% to hold an argument's detokenized string and its uppercased form,
% for comparison.
%    \begin{macrocode}
\tl_new:N \l_lsm_str_tl
\tl_new:N \l_lsm_upper_tl
%    \end{macrocode}
% \end{variable}
%
% \begin{macro}[no-user-doc]{\lsm_if_arbitrary_symbol:nTF}
% Shared classification test behind \cs{catName}'s and
% \cs{morphName}\slash \cs{morphSeqName}'s auto-detection
% (\pagecrefn{sec:style-impl}): is \texttt{\#1} an \emph{arbitrary symbol}
% (an ordinary bare-symbol name, however decorated) as opposed to a
% genuine multi-letter \emph{word}? (\pagecrefn{sec:cat-symbols}):
% true if the argument contains \texttt{\^{}} or \texttt{\_{}} (a
% sub/superscript decoration, regardless of what the decoration itself
% contains), or is a single token, or consists entirely of upper-case
% letters; false otherwise. A hand-written manual conditional (not
% declared via \cs{prg_new_conditional:Npnn}, which would need a
% \texttt{p} outcome this test never requires) since its branches are
% used exactly like an ordinary \texttt{TF} conditional's.
%    \begin{macrocode}
\cs_new_protected:Npn \lsm_if_arbitrary_symbol:nTF #1#2#3
  {
    \str_if_in:nnTF {#1} { ^ }
      {#2}
      {
        \str_if_in:nnTF {#1} { _ }
          {#2}
          {
            \tl_if_single:nTF {#1}
              {#2}
              {
                \tl_set:Nx \l_lsm_str_tl { \tl_to_str:n {#1} }
                \tl_set:Nx \l_lsm_upper_tl
                  { \str_uppercase:f { \tl_to_str:n {#1} } }
                \tl_if_eq:NNTF \l_lsm_str_tl \l_lsm_upper_tl
                  {#2}
                  {#3}
              }
          }
      }
  }
%    \end{macrocode}
% \end{macro}
%
% \subsection{Public commands}
%
% Every user-level command is declared once, unconditionally, here.
% None of them change between styles: each is a fixed, thin wrapper
% around an internal presentation-hook function of the same shape,
% and it is only those internal functions that a style reassigns
% (\pagecrefn{sec:style-switching} above; \pagecrefn{sec:style-impl} below).
% \begin{macro}{\semName, \catName, \catSeqName, \objName, \morphName,
%   \morphSeqName, \morphMono, \morphEpi, \morphCompose,
%   \morphComposeHead, \morphComposeTail, \morphId, \morphSeqId,
%   \morphIncl, \morphSeqIncl, \morphInclArrow, \chartName, \atlasName,
%   \coordName}
% \cs{morphIncl} and \cs{morphSeqIncl} split their optional comma list
% into the included object and the containing object before calling
% their hooks; more than two items is an error.
%    \begin{macrocode}
\NewDocumentCommand \semName { m } { \lsm_sem_name:n {#1} }
\NewDocumentCommand \catName { m } { \lsm_cat_name:n {#1} }
\NewDocumentCommand \catSeqName { m } { \lsm_cat_seqname:n {#1} }
\NewDocumentCommand \objName { m } { \lsm_obj_name:n {#1} }
\NewDocumentCommand \morphName { m } { \lsm_morph_name:n {#1} }
\NewDocumentCommand \morphSeqName { m } { \lsm_morph_seqname:n {#1} }
\NewDocumentCommand \morphMono { } { \lsm_morph_mono: }
\NewDocumentCommand \morphEpi { } { \lsm_morph_epi: }
\NewDocumentCommand \morphCompose { O{} } { \lsm_morph_compose:n {#1} }
\NewDocumentCommand \morphComposeHead { O{} }
  { \lsm_morph_compose_head:n {#1} }
\NewDocumentCommand \morphComposeTail { O{} }
  { \lsm_morph_compose_tail:n {#1} }
\NewDocumentCommand \morphId { O{} } { \lsm_morph_id:n {#1} }
\NewDocumentCommand \morphSeqId { O{} } { \lsm_morph_seqid:n {#1} }
\NewDocumentCommand \morphIncl { O{} }
  { \lsm_incl_split:Nnn \lsm_morph_incl:nn { morphIncl } {#1} }
\NewDocumentCommand \morphSeqIncl { O{} }
  { \lsm_incl_split:Nnn \lsm_morph_seqincl:nn { morphSeqIncl } {#1} }
\NewDocumentCommand \morphInclArrow { } { \lsm_morph_incl_arrow: }

\clist_new:N \l_lsm_incl_tmp_clist
\msg_new:nnn { shmuelsemtex } { incl-items }
  { More~than~two~items~for~\c_backslash_str #1:~[#2] }
\cs_new_protected:Npn \lsm_incl_split:Nnn #1 #2 #3
  {
    \clist_set:Nn \l_lsm_incl_tmp_clist {#3}
    \int_case:nnF { \clist_count:N \l_lsm_incl_tmp_clist }
      {
        { 0 } { #1 { } { } }
        { 1 }
          {
            \exp_args:Ne #1
              { \clist_item:Nn \l_lsm_incl_tmp_clist { 1 } } { }
          }
        { 2 }
          {
            \exp_args:Nee #1
              { \clist_item:Nn \l_lsm_incl_tmp_clist { 1 } }
              { \clist_item:Nn \l_lsm_incl_tmp_clist { 2 } }
          }
      }
      {
        \msg_error:nnnn { shmuelsemtex } { incl-items } {#2} {#3}
      }
  }
\NewDocumentCommand \chartName { m } { \lsm_chart_name:n {#1} }
\NewDocumentCommand \atlasName { m } { \lsm_atlas_name:n {#1} }
\NewDocumentCommand \coordName { m } { \lsm_coord_name:n {#1} }
%    \end{macrocode}
% \end{macro}
%
% \subsection{Style implementations}
% \label{sec:style-impl}
%
% Each style is one function that reassigns \emph{every} presentation
% hook in the package via \cs{cs_set_protected:Npn} (not
% \cs{cs_new_protected:Npn}: styles overwrite, they do not merely
% declare, and no pre-existence check is wanted since more than one
% style will assign the same hook name). A style is applied by
% calling its function, never by patching individual hooks.
%
% \begin{macro}[no-user-doc]{\lsm_style_default:}
% The \texttt{default} style: the presentation choices justified,
% macro by macro, in the documentation above.
%    \begin{macrocode}
\cs_new_protected:Npn \lsm_style_default:
  {
    \cs_set_protected:Npn \lsm_sem_name:n ##1 { ##1 }
    \cs_set_protected:Npn \lsm_cat_name:n ##1
      {
        \lsm_if_arbitrary_symbol:nTF {##1}
          { \mathscr{##1} }
          { \mathbf{##1} }
      }
    \cs_set_protected:Npn \lsm_cat_seqname:n ##1 { \bm{\mathscr{##1}} }
    \cs_set_protected:Npn \lsm_obj_name:n ##1 { ##1 }
    \cs_set_protected:Npn \lsm_morph_name:n ##1
      {
        \lsm_if_arbitrary_symbol:nTF {##1}
          { \mathit{##1} }
          { \mathrm{##1} }
      }
    \cs_set_protected:Npn \lsm_morph_seqname:n ##1
      {
        \lsm_if_arbitrary_symbol:nTF {##1}
          { \bm{\mathit{##1}} }
          { \bm{\mathrm{##1}} }
      }
    \cs_set_protected:Npn \lsm_morph_mono: { \rightarrowtail }
    \cs_set_protected:Npn \lsm_morph_epi: { \twoheadrightarrow }
    \cs_set_protected:Npn \lsm_morph_compose:n ##1
      {
        \tl_if_empty:nTF {##1}
          { \mathbin{\circ} }
          { \mathbin{\overset{##1}{\circ}} }
      }
    \cs_set_protected:Npn \lsm_morph_compose_head:n ##1
      {
        \tl_if_empty:nTF {##1}
          { \mathbin{\odot} }
          { \mathbin{\overset{##1}{\odot}} }
      }
    \cs_set_protected:Npn \lsm_morph_compose_tail:n ##1
      {
        \tl_if_empty:nTF {##1}
          { \mathbin{\cdot} }
          { \mathbin{\overset{##1}{\cdot}} }
      }
    \lsm_style_id_incl:
    \cs_set_protected:Npn \lsm_chart_name:n ##1 { ##1 }
    \cs_set_protected:Npn \lsm_atlas_name:n ##1 { ##1 }
    \cs_set_protected:Npn \lsm_coord_name:n ##1 { ##1 }
  }
%    \end{macrocode}
% \end{macro}
%
% \begin{macro}[no-user-doc]{\lsm_style_id_incl:}
% The identity and inclusion hooks, shared by both built-in styles:
% upright Id, bold for a sequence, with the identity's object as a
% subscript and an inclusion's included and containing objects as
% superscript and subscript; the inclusion arrow is a hook arrow.
% Identity and inclusion are separate hooks even where, as here, they
% render alike, so that a style can tell them apart.
%    \begin{macrocode}
\cs_new_protected:Npn \lsm_id_incl:nnn #1 #2 #3
  {
    \operatorname{#1}
    \tl_if_empty:nF {#2} { \sp {#2} }
    \tl_if_empty:nF {#3} { \sb {#3} }
  }
\cs_new_protected:Npn \lsm_style_id_incl:
  {
    \cs_set_protected:Npn \lsm_morph_id:n ##1
      { \lsm_id_incl:nnn { Id } { } {##1} }
    \cs_set_protected:Npn \lsm_morph_seqid:n ##1
      { \lsm_id_incl:nnn { \mathbf{Id} } { } {##1} }
    \cs_set_protected:Npn \lsm_morph_incl:nn ##1 ##2
      { \lsm_id_incl:nnn { Id } {##1} {##2} }
    \cs_set_protected:Npn \lsm_morph_seqincl:nn ##1 ##2
      { \lsm_id_incl:nnn { \mathbf{Id} } {##1} {##2} }
    \cs_set_protected:Npn \lsm_morph_incl_arrow: { \hookrightarrow }
  }
%    \end{macrocode}
% \end{macro}
%
% \begin{macro}[no-user-doc]{\lsm_style_plain:}
% The \texttt{plain} style: every \texttt{*Name} family macro
% unstyled; mono/epi/compose glyphs unchanged from \texttt{default}
% (see the documentation above for why). Demonstrates that a style is
% a real, independent, complete bundle, not a variant of
% \cs{lsm_style_default:}.
%    \begin{macrocode}
\cs_new_protected:Npn \lsm_style_plain:
  {
    \cs_set_protected:Npn \lsm_sem_name:n ##1 { ##1 }
    \cs_set_protected:Npn \lsm_cat_name:n ##1 { ##1 }
    \cs_set_protected:Npn \lsm_cat_seqname:n ##1 { ##1 }
    \cs_set_protected:Npn \lsm_obj_name:n ##1 { ##1 }
    \cs_set_protected:Npn \lsm_morph_name:n ##1 { ##1 }
    \cs_set_protected:Npn \lsm_morph_seqname:n ##1 { ##1 }
    \cs_set_protected:Npn \lsm_morph_mono: { \rightarrowtail }
    \cs_set_protected:Npn \lsm_morph_epi: { \twoheadrightarrow }
    \cs_set_protected:Npn \lsm_morph_compose:n ##1
      {
        \tl_if_empty:nTF {##1}
          { \mathbin{\circ} }
          { \mathbin{\overset{##1}{\circ}} }
      }
    \cs_set_protected:Npn \lsm_morph_compose_head:n ##1
      {
        \tl_if_empty:nTF {##1}
          { \mathbin{\odot} }
          { \mathbin{\overset{##1}{\odot}} }
      }
    \cs_set_protected:Npn \lsm_morph_compose_tail:n ##1
      {
        \tl_if_empty:nTF {##1}
          { \mathbin{\cdot} }
          { \mathbin{\overset{##1}{\cdot}} }
      }
    \lsm_style_id_incl:
    \cs_set_protected:Npn \lsm_chart_name:n ##1 { ##1 }
    \cs_set_protected:Npn \lsm_atlas_name:n ##1 { ##1 }
    \cs_set_protected:Npn \lsm_coord_name:n ##1 { ##1 }
  }
%    \end{macrocode}
% \end{macro}
%
% \subsection{Package options and runtime setup}
%
% \begin{macro}{style}
% A single \texttt{.choice:} key in the \texttt{lsm} key family,
% defaulting to \texttt{default} via \texttt{.initial:n} (which runs
% immediately at definition time, so the package works correctly with
% no options given at all), then re-applied if the author gives an
% explicit \texttt{style=...} package option, or calls
% \cs{setupsemanticmarkup} (below) at any later point.
%    \begin{macrocode}
\keys_define:nn { lsm }
  {
    style .choice: ,
    style / default .code:n = { \lsm_style_default: } ,
    style / plain .code:n = { \lsm_style_plain: } ,
    style .initial:n = default ,
  }
\ProcessKeysPackageOptions { lsm }
%    \end{macrocode}
% \end{macro}
%
% \begin{macro}{\setupsemanticmarkup}
% Runtime counterpart to the package option above, in the manner of
% \cs{setupquant}\slash \cs{setupset}: a thin wrapper around
% \cs{keys_set:nn} on the same key family, callable anywhere in the
% document body, not only on the \cs{usepackage} line.
%    \begin{macrocode}
\NewDocumentCommand \setupsemanticmarkup { m }
  { \keys_set:nn { lsm } {#1} }
%    \end{macrocode}
% \end{macro}
%
% \subsection{Quantifiers, sets, and sequence names}
% \label{sec:quant-set-impl}
%
% The public interface of these commands (\cs{equant}, \cs{uquant},
% \cs{set}, \cs{setupquant}, \cs{setupset}, \cs{union},
% \cs{intersection}, \cs{seqname}) is documented above;
% \cs{setupquant} and \cs{setupset} are the model for
% \cs{setupsemanticmarkup} (see the discussion in
% \pagecrefn{sec:style-impl}). Each macro family has its own scoped,
% uniquely named scratch variables, so one family nested inside
% another's argument (e.g.\ a \cs{union} inside a \cs{set} predicate)
% is safe.
%
% \begin{variable}[no-user-doc]{\l_lsm_quant_tmp_clist, \l_lsm_quant_tmp_fp,
%   \l_lsm_set_tmp_clist, \l_lsm_set_tmp_tl, \l_lsm_set_star_bool,
%   \l_lsm_unint_tmp_clist, \l_lsm_unint_tmp_fp}
% Each macro family below gets its own uniquely-named scratch
% variables, declared together here rather than reusing a shared
% \texttt{tmpa}-style name across families -- the fix for the latent
% nesting bug described above.
%    \begin{macrocode}
\clist_new:N \l_lsm_quant_tmp_clist
\fp_new:N    \l_lsm_quant_tmp_fp
\clist_new:N \l_lsm_set_tmp_clist
\tl_new:N    \l_lsm_set_tmp_tl
\bool_new:N  \l_lsm_set_star_bool
\clist_new:N \l_lsm_unint_tmp_clist
\fp_new:N    \l_lsm_unint_tmp_fp
%    \end{macrocode}
% \end{variable}
%
% \begin{macro}{\equant, \uquant, \setupquant}
% Existential/universal quantifier rendering, with a runtime style
% switch (parenthesization and subscript-vs-inline placement of the
% bound variables) via \cs{setupquant}, independent of
% \cs{setupsemanticmarkup} above; it is
% kept as its own key family rather than folded into
% \texttt{lsm}'s, since it predates and is unrelated to the
% default/plain style choice.
%    \begin{macrocode}
\int_new:N \g_lsm_quant_parens_int
\int_new:N \g_lsm_quant_subscr_int
\tl_new:N \g_lsm_quant_sep_tl

\NewDocumentCommand{\equant}{mm}
  { \lsm_quant:nnn {\exists} {#1} {#2} }

\NewDocumentCommand{\uquant}{mm}
  { \lsm_quant:nnn {\forall} {#1} {#2} }

\NewDocumentCommand \setupquant {m}
  { \keys_set:nn {lsm-quant} {#1} }

\keys_define:nn {lsm-quant}
 {
  subscript   .choices:nn =
    {
      {none, stacked, multiple}
      { \int_gset:Nn \g_lsm_quant_subscr_int {\l_keys_choice_int - 1} }
    },
  subscript   .default:n = multiple,
  subscript   .initial:n = none,
  parentheses .choices:nn =
    {
      {none, single, multiple}
      { \int_gset:Nn \g_lsm_quant_parens_int {\l_keys_choice_int - 1} }
    },
  parentheses .default:n = multiple,
  parentheses .initial:n = none,
  separator   .tl_set:N  = \g_lsm_quant_sep_tl,
  separator   .default:n = {.},
  separator   .initial:n = {},
 }

\cs_new_protected:Npn \lsm_quant:nnn #1 #2 #3
  {
    \clist_set:Nn \l_lsm_quant_tmp_clist {#2}
    \int_case:nn {\g_lsm_quant_subscr_int}
      {
        {0}
        {
          \int_compare:nT {\g_lsm_quant_parens_int > 0} {\left (}
          #1
          \clist_use:Nn \l_lsm_quant_tmp_clist {\g_lsm_quant_sep_tl}
          \int_compare:nT {\g_lsm_quant_parens_int > 0} {\right )}
          \g_lsm_quant_sep_tl #3
        }
        {1}
        {
          \fp_set:Nn \l_lsm_quant_tmp_fp
            {ceil{\clist_count:N \l_lsm_quant_tmp_clist - 1} * .1 + 1}
          \int_compare:nT {\g_lsm_quant_parens_int > 0} {\left (}
          \scaleobj{\fp_to_decimal:N \l_lsm_quant_tmp_fp}{#1} \sb
            { \substack { \clist_use:Nn \l_lsm_quant_tmp_clist { \\ } } }
          \int_compare:nT {\g_lsm_quant_parens_int > 0} {\right )}
          #3
        }
        {2}
        {
          \int_compare:nT {\g_lsm_quant_parens_int = 1} {\left (}
          \clist_map_inline:Nn \l_lsm_quant_tmp_clist
            {
              \int_compare:nT {\g_lsm_quant_parens_int = 2} {\left (}
              \scaleobj{1.2}{#1} \sb {##1}
              \int_compare:nT {\g_lsm_quant_parens_int = 2} {\right )}
            }
          \int_compare:nT {\g_lsm_quant_parens_int = 1} {\right )}
          #3
        }
      }
  }
%    \end{macrocode}
% \end{macro}
%
% \begin{macro}{\set, \setupset}
% Set-builder notation, \cs{set}\marg{elements}\oarg{propositions},
% with an optional star after the arguments
% (\cs{set}\marg{elements}\oarg{propositions}\texttt{*}) that splits the
% line between elements and propositions -- useful inside a
% \env{multline} environment, since \LaTeX{} cannot otherwise break a
% line inside a macro argument. \cs{setupset}'s \texttt{separator} key
% sets the symbol between elements and propositions; its
% \texttt{subscript} key applies to \cs{union} and \cs{intersection}
% (below).
%    \begin{macrocode}
\tl_new:N \g_lsm_set_sep_tl
\int_new:N \g_lsm_set_subscr_int

\NewDocumentCommand{\set}{mos}
  {
    \IfBooleanTF {#3}
      { \bool_set_true:N  \l_lsm_set_star_bool }
      { \bool_set_false:N \l_lsm_set_star_bool }
    \lsm_set_of:nnn {#1} {#2} {\l_lsm_set_star_bool}
  }

\NewDocumentCommand \setupset {m}
  { \keys_set:nn {lsm-set} {#1} }

\keys_define:nn {lsm-set}
  {
    separator .tl_set:N  = \g_lsm_set_sep_tl,
    separator .initial:n = {\mid},
    subscript .choices:nn =
      {
        {stacked, multiple}
        { \int_gset:Nn \g_lsm_set_subscr_int {\l_keys_choice_int - 1} }
      },
    subscript .initial:n = stacked,
  }

\cs_new_protected:Npn \lsm_set_of:nnn #1 #2 #3
  {
    \IfValueTF {#2}
      {
        \clist_set:Nn \l_lsm_set_tmp_clist {#2}
        \tl_set:Nx \l_lsm_set_tmp_tl
          { \clist_use:Nn \l_lsm_set_tmp_clist {\land} }
        \bool_if:nTF {#3}
          {
            \scalerel*{\{}{#1 \l_lsm_set_tmp_tl}
            #1
            \clist_set:Nn \l_lsm_set_tmp_clist {#2}
            \scalerel*{\mathbin{\g_lsm_set_sep_tl}}{#1 \l_lsm_set_tmp_tl}
            \\
            \clist_set:Nn \l_lsm_set_tmp_clist {#2}
            \l_lsm_set_tmp_tl
            \clist_set:Nn \l_lsm_set_tmp_clist {#2}
            \scalerel*{\}}{#1 \l_lsm_set_tmp_tl}
          }
          {
            \left \{
            #1
            \clist_set:Nn \l_lsm_set_tmp_clist {#2}
            \scalerel*{\mathbin{\g_lsm_set_sep_tl}}{#1 \l_lsm_set_tmp_tl}
            \l_lsm_set_tmp_tl
            \right \}
          }
      }
      { \left \{ #1 \right \} }
  }
%    \end{macrocode}
% \end{macro}
%
% \begin{macro}{\union, \intersection}
% \cs{union}\oarg{indices}\marg{set} and
% \cs{intersection}\oarg{indices}\marg{set} typeset indexed union and intersection; the index list's subscript
% placement (stacked under the operator, or inline comma-separated) is
% controlled by \cs{setupset}'s \texttt{subscript} key above, shared
% with \cs{set} since both are the same family of indexed-operator
% notation.
%    \begin{macrocode}
\NewDocumentCommand{\intersection}{om}
  { \lsm_unint_of:nnn \bigcap {#1} {#2} }

\NewDocumentCommand{\union}{om}
  { \lsm_unint_of:nnn \bigcup {#1} {#2} }

\cs_new_protected:Npn \lsm_unint_of:nnn #1 #2 #3
  {
    \IfValueTF {#2}
      {
        \clist_set:Nn \l_lsm_unint_tmp_clist {#2}
        \int_case:nn {\g_lsm_set_subscr_int}
          {
            {0}
            {
              \fp_set:Nn \l_lsm_unint_tmp_fp
                {ceil{\clist_count:N \l_lsm_unint_tmp_clist - 1} * .1 + 1}
              \scaleobj{\fp_to_decimal:N \l_lsm_unint_tmp_fp}{#1} \sb
                { \substack { \clist_use:Nn \l_lsm_unint_tmp_clist { \\ } } }
              #3
            }
            {1}
            {
              #1 \sb {\clist_use:Nn \l_lsm_unint_tmp_clist {,}}
              #3
            }
          }
      }
      { #1 #3 }
  }
%    \end{macrocode}
% \end{macro}
%
% \begin{macro}{\seqname}
% Render a sequence/set/tuple name: bold upright for a multi-character
% name, bold italic for a single character -- the same
% length-based split \cs{catName} uses, kept as its own macro since it
% is semantically a different kind of object (a sequence, not a
% category) even though the rendering rule happens to coincide.
%    \begin{macrocode}
\NewDocumentCommand{\seqname}{m}
  { \lsm_seqname:n {#1} }

\cs_new_protected:Npn \lsm_seqname:n #1
  {
    \int_compare:nTF {\tl_count:n {#1} > 1}
      { {\mathbf{#1}} }
      {
        \cs_if_exist:NTF \mathbfit
          { {\mathbfit{#1}} }
          { {\bm{\mathit{#1}}} }
      }
  }
%    \end{macrocode}
% \end{macro}
%
%</package>
% \end{implementation}
%
% \PrintIndex
